\subsection{COMPACT METRIC SPACES; COMPACT METRIZABLE SPACES}

The criterion for precompactness of a uniform space (Chapter II, § 4, no. 2, Theorem 3) gives rise to the following proposition for metric spaces:

PROPOSITION 14. A metric space X is precompact if and only if, for each \(\varepsilon > 0\) , there is a finite covering of X by sets of diameter \(\leqslant \varepsilon\) .

If we adjoin the hypothesis that X is complete we have a criterion for compactness of metric spaces.

From Proposition 14 we get a topological criterion for compactness, applicable to metrizable spaces:

PROPOSITION 15. A metrizable topological space X is compact if and only if every infinite sequence of points of X has a cluster point in X.

Axiom (C) of Chapter I, § 9, no. 1 shows that the condition is necessary. To show sufficiency, let \(d\) be a metric compatible with the topology of \(\mathbf{X}\) . We show first that the metric space \(\mathbf{X}\) so defined is complete: every Cauchy sequence in \(\mathbf{X}\) has a cluster point and is therefore convergent (Chapter II, § 3, no. 2, Proposition 5, Corollary 2); hence \(\mathbf{X}\) is complete, by Proposition 9. Next we shall show that \(\mathbf{X}\) is precompact; if this were not so, then by Proposition 14 there would exist a real number \(\alpha > 0\) such that \(\mathbf{X}\) could not be covered by any finite number of subsets of \(\mathbf{X}\) of diameter \(\leqslant \alpha\) . We could then define by induction on \(n\) an infinite sequence \((x_n)\) of points of \(\mathbf{X}\) by the condition \(d(x_p, x_n) > \frac{1}{2} \alpha\) for all \(p < n\) ; and such a sequence can have no cluster point, since every ball of radius \(< \frac{1}{2} \alpha\) contains at most one point of the sequence.

COROLLARY. A subset A of a metrizable topological space X is relatively compact if and only if every infinite sequence of points of A has a cluster point in X.

Let \(d\) be a metric compatible with the topology of \(\mathbf{X}\) . We shall show that the space \(\overline{\mathbf{A}}\) is compact, by applying the criterion of Proposition 15. Let \((x_{n})\) be a sequence of points of \(\overline{\mathbf{A}}\) ; then for each index \(n\) there exists \(y_{n} \in \mathbf{A}\) such that \(d(x_{n}, y_{n}) < 1/n\) ; the sequence \((y_{n})\) has, by hypothesis, a cluster point \(a \in \mathbf{X}\) , and \(a\) is also a cluster point of the sequence \((x_{n})\) , for if \(y_{m}\) lies in the ball with centre \(a\) and radius \(1/n\) , for some \(m > n\) , then \(x_{m}\) lies in the ball with centre \(a\) and radius \(2/n\) .

It should be remarked that Proposition 15 is not a consequence of the existence of a countable fundamental system of neighbourhoods at each point of X; there are examples of non-metrizable, non-compact spaces in which every point has a countable fundamental system of neighbourhoods and every infinite sequence of points has a cluster point (Exercise 15).

PROPOSITION 16. A compact space X is metrizable if and only if its topology has a countable base.

The condition is necessary. By Proposition 14, for each integer \(n \geqslant 1\) there is a finite subset \(A_n\) of \(X\) such that the distance of \(A_n\) from every point of \(X\) is \(\leqslant 1/n\) ; the countable set \(A = \bigcup_{n} A_n\) is therefore dense in \(X\) , and the result follows from Proposition 12 of no. 8.

The condition is sufficient. Let \((\mathbf{U}_{n})\) be a countable base of the topology of X. Every neighbourhood of a point of the diagonal \(\Delta\) of \(X \times X\) therefore contains a neighbourhood of the form \(U_{n} \times U_{n}\) ; applying the Borel-Lebesgue axiom to the compact subset \(\Delta\) of \(X \times X\) , it follows that every neighbourhood of \(\Delta\) contains a finite union of sets of the form \(U_{n} \times U_{n}\) , which is a neighbourhood of \(\Delta\) . Hence the neighbourhoods of \(\Delta\) which are finite unions of sets of the form \(U_{n} \times U_{n}\) form a fundamental system of entourages of the uniformity of X (Chapter II, § 4, no. 1, Theorem 1), and the result therefore follows from Theorem 1 of no. 4.

COROLLARY. Let X be a locally compact space and let X' be the compact space obtained by adjoining a point at infinity ω to X (Chapter I, § 9, no. 8). Then the following statements are equivalent:

\begin{enumerate}
\def\labelenumi{\alph{enumi})}
\item
  The topology of \(\mathbf{X}\) has a countable base.
\item
  \(\mathbf{X}'\) is metrizable.
\item
  X is metrizable and \(\sigma\) -compact.
\end{enumerate}

a) \(\Longrightarrow\) b): Let \((\mathbf{U}_n)\) be a countable base of the topology of \(\mathbf{X}\) . Each neighbourhood of a point \(x \in \mathbf{X}\) contains a compact neighbourhood of \(x\) , which in turn contains a neighbourhood of \(x\) equal to some \(\mathbf{U}_n\) . Hence the relatively compact \(\mathbf{U}_n\) form a base of the topology of \(\mathbf{X}\) , and we may therefore suppose that all the \(\mathbf{U}_n\) are relatively compact. \(\mathbf{X}\) is therefore a countable union of compact sets \(\overline{\mathbf{U}}_n\) , i.e.~it is \(\sigma\) -compact; this implies that in \(\mathbf{X}'\) the point \(\omega\) has a countable fundamental system \((\mathbf{V}_n)\) of open neighbourhoods (Chapter I, § 9, no. 9, Proposition 15, Corollary 2). Hence each neighbourhood of a point \(y \in \mathbf{X}'\) contains either one of the \(\mathbf{U}_n\) or one of the \(\mathbf{V}_n\) , which is a neighbourhood of \(y\) , and so the \(U_{n}\) and the \(V_{n}\) form a countable base of the topology of \(X'\) . Hence \(X'\) is metrizable, by Proposition 16.

b) \(\Longrightarrow\) c): If \(X'\) is metrizable, then \(\omega\) has a countable fundamental system of neighbourhoods, and therefore \(X\) is \(\sigma\) -compact by Chapter I, § 9, no. 9, Proposition 15, Corollary 2.

c) \(\Longrightarrow a)\) : By hypothesis, there is an increasing sequence \((\mathbf{V}_n)\) of relatively compact open sets which cover \(\mathbf{X}\) and are such that \(\overline{\mathbf{V}}_n \subset \mathbf{V}_{n+1}\) (Chapter I, § 9, no. 9, Proposition 15). The subspace \(\overline{\mathbf{V}}_n\) is compact and metrizable and therefore has a countable base (Proposition 16), and therefore so does \(\mathbf{V}_n\) . Let \((\mathbf{U}_{mn})_{m \geqslant 1}\) be a base of the topology of \(\mathbf{V}_n\) . For each \(x \in \mathbf{X}\) and each neighbourhood \(W\) of \(x\) , there exists \(n\) such that \(x \in \mathbf{V}_n\) , hence there exists \(m\) such that \(x \in \mathbf{U}_{mn} \subset \mathbf{V}_n \cap W\) . Hence the sets \(\mathbf{U}_{mn}\) ( \(m \geqslant 1\) , \(n \geqslant 1\) ) form a base of the topology of \(\mathbf{X}\) .
% semantic-fragments: legacy-input-seam
\relax{}
